% Benchmark results — core-guided vs standard MUS
% Generated from bench-results/ directory, 2026-04-30

% =============================================================
%  TABLE
% =============================================================

\begin{table*}[t]
\centering
\caption{Comparison of two MUS-finding strategies across six puzzle
  types (mean over 20 instances per row).  \emph{Solve time} is
  wall-clock time for the solve phase only (excluding Conjure
  compilation and parse setup).  \emph{Max MUS} is the largest MUS
  returned in any step of the explanation.  The ``OA'' rows use
  \texttt{--only-assign}, restricting deductions to cell assignments.}
\label{tab:mus-comparison}
\small
\begin{tabular}{@{}ll rr rr@{}}
\toprule
 & & \multicolumn{2}{c}{Solve time (s)}
 & \multicolumn{2}{c}{Max MUS size} \\
\cmidrule(lr){3-4} \cmidrule(lr){5-6}
Puzzle & Difficulty
 & \textsc{mus} & \textsc{core}
 & \textsc{mus} & \textsc{core} \\
\midrule
% --- Binairo ---
Binairo
  & 6$\times$6 easy    & 0.01 & 0.01  &  1.0 &  1.0 \\
  & 6$\times$6 hard    & 0.01 & 0.01  &  2.2 &  3.1 \\
  & 10$\times$10 easy  & 0.08 & 0.08  &  1.0 &  1.0 \\
  & 10$\times$10 hard  & 0.12 & 0.09  &  2.0 &  3.8 \\
  & 20$\times$20 easy  & 2.50 & 2.22  &  1.1 &  1.1 \\
  & 20$\times$20 hard  & 3.71 & 2.58  &  2.0 &  6.0 \\
\addlinespace
% --- Sudoku ---
Sudoku
  & basic              & 1.09 & 1.04  &  1.0 &  1.2 \\
  & easy               & 1.18 & 1.08  &  1.0 &  1.1 \\
  & intermediate       & 1.35 & 1.05  &  2.5 &  4.0 \\
  & advanced           & 1.83 & 1.09  &  2.5 &  7.0 \\
  & extreme            & 2.37 & 1.04  &  4.6 &  7.5 \\
  & evil               & 5.99 & 1.15  &  6.0 & 16.3 \\
\addlinespace
% --- Futoshiki ---
Futoshiki
  & 4$\times$4 easy    & 0.01 & 0.01  &  1.0 &  1.0 \\
  & 5$\times$5 easy    & 0.03 & 0.03  &  1.0 &  1.0 \\
  & 5$\times$5 normal  & 0.04 & 0.03  &  2.1 &  4.0 \\
  & 5$\times$5 hard    & 0.08 & 0.03  &  3.9 &  6.3 \\
  & 9$\times$9 easy    & 1.94 & 1.22  &  1.0 &  1.0 \\
  & 9$\times$9 normal  & 3.01 & 1.27  &  2.8 & 10.4 \\
  & 9$\times$9 hard    & 9.52 & 1.37  &  5.5 & 14.2 \\
\addlinespace
% --- Akari ---
Akari
  & 7$\times$7 easy    & 0.01 & 0.01  &  1.0 &  1.0 \\
  & 7$\times$7 normal  & 0.01 & 0.01  &  3.0 &  3.6 \\
  & 7$\times$7 hard    & 0.02 & 0.01  &  4.7 &  5.2 \\
  & 10$\times$10 easy  & 0.03 & 0.03  &  1.0 &  1.0 \\
  & 10$\times$10 normal& 0.05 & 0.04  &  3.0 &  3.5 \\
  & 10$\times$10 hard  & 0.07 & 0.04  &  5.5 &  8.4 \\
  & 14$\times$14 easy  & 0.14 & 0.12  &  1.0 &  1.1 \\
  & 14$\times$14 normal& 0.23 & 0.14  &  3.0 &  3.9 \\
  & 14$\times$14 hard  & 0.33 & 0.14  &  4.9 &  6.8 \\
\addlinespace
% --- Minesweeper ---
Minesweeper
  & 5$\times$5 easy    & 0.00 & 0.00  &  1.0 &  1.0 \\
  & 5$\times$5 hard    & 0.00 & 0.00  &  1.8 &  1.9 \\
  & 10$\times$10 easy  & 0.02 & 0.02  &  1.0 &  1.1 \\
  & 10$\times$10 hard  & 0.02 & 0.02  &  2.0 &  3.0 \\
  & 20$\times$20 easy  & 0.25 & 0.22  &  1.0 &  1.1 \\
  & 20$\times$20 hard  & 0.46 & 0.28  &  2.0 &  3.6 \\
\addlinespace
% --- Tents ---
Tents
  & 6$\times$6 easy    & 0.27 & 0.29  &  1.1 &  1.2 \\
  & 6$\times$6 hard    & 0.29 & 0.28  &  2.5 &  3.8 \\
  & 10$\times$10 easy  &25.4  &26.0   &  1.4 &  1.9 \\
  & 10$\times$10 hard  &25.5  &25.9   &  2.5 &  5.8 \\
\midrule
% --- Futoshiki OA ---
Futoshiki (OA)
  & 4$\times$4 easy    & 0.01 & 0.00  &  3.2 &  3.2 \\
  & 5$\times$5 easy    & 0.04 & 0.01  &  4.8 &  4.9 \\
  & 5$\times$5 normal  & 0.04 & 0.01  &  7.8 &  9.8 \\
  & 5$\times$5 hard    & 0.05 & 0.02  & 10.9 & 13.9 \\
  & 9$\times$9 easy    & 3.21 & 0.52  & 11.4 & 11.9 \\
  & 9$\times$9 normal  & 3.63 & 0.54  & 18.5 & 46.2 \\
  & 9$\times$9 hard    & 4.00 & 0.54  & 25.9 & 60.6 \\
\addlinespace
% --- Tents OA ---
Tents (OA)
  & 6$\times$6 easy    & 0.06 & 0.05  &  2.8 &  4.5 \\
  & 6$\times$6 hard    & 0.11 & 0.05  &  4.9 &  7.4 \\
  & 10$\times$10 easy  & 5.47 & 2.02  &  4.8 &16.5  \\
  & 10$\times$10 hard  & 9.63 & 2.18  &  5.8 &19.9  \\
\bottomrule
\end{tabular}
\end{table*}


% =============================================================
%  EXPERIMENTAL SETUP
% =============================================================

\subsection{Experimental setup}

We compare two strategies for computing the constraint-based
explanations produced by Demystify:

\begin{description}
\item[\textsc{mus}] Standard deletion-based MUS extraction.  For each
  deduction, we enumerate all constraints that could be relevant, then
  iteratively remove constraints one at a time, keeping only those
  whose removal makes the remaining set satisfiable.

\item[\textsc{core}] Core-guided extraction.  We ask the SAT solver
  for an unsatisfiable core directly and return it without further
  minimisation.  Cores are not guaranteed to be minimal.
\end{description}

We also tested a hybrid strategy that primes the deletion-based
minimisation with a SAT core (rather than starting from the full
constraint set), but it neither reduced solve time nor MUS size
beyond \textsc{mus} on any puzzle family above 0.5\,s --- aggregate
solve time differed by under~2\% and per-row max-MUS sizes were
indistinguishable.  We omit it from the table.

Both strategies use the same underlying SAT solver (Glucose~4,
via the \texttt{rustsat-glucose} bindings) with a per-call conflict
limit of 100, ramped up adaptively when too many calls return
inconclusive (see Section~\ref{sec:enhancements}).  Explanations are
generated step-by-step: at each step the solver identifies all forced
deductions, computes a MUS (or core) justifying each one, then fixes
the deduced values before proceeding to the next step.  There is no
step limit.

The benchmark corpus comprises 20 instances at each difficulty level
for six puzzle types: Binairo, Sudoku, Futoshiki, Akari (Light Up),
Minesweeper, and Tents.  Instances were sourced from online puzzle
generators (puzzle-binairo.com, puzzle-sudoku.com, etc.).  Each puzzle
type is modelled in Essence Prime and compiled to CNF via Conjure and
Savile Row.  We also test Futoshiki and Tents in ``only-assign'' (OA)
mode, where deductions are restricted to cell assignments (ignoring
domain eliminations), producing fewer but harder-to-justify steps.

All experiments ran on an Apple M4 Mac mini (10 cores, 16\,GB RAM)
running macOS~26.4, compiled with Rust~1.95.0 in release mode.
Timings are wall-clock time for the solve phase only, excluding the
Conjure compilation and parse setup (which are identical across
strategies).  Each value in the table is the mean over 20 instances.


% =============================================================
%  ENHANCEMENTS
% =============================================================

\subsection{Solver enhancements}
\label{sec:enhancements}

Two changes to the solve loop materially affect the numbers in
Table~\ref{tab:mus-comparison}.

\paragraph{Expanding a MUS to all literals it justifies.}

A MUS is computed to explain a single forced deduction, but the same
constraint set typically rules out other literals as well.  Once a MUS
$M$ has been minimised, we ask the SAT solver --- under the
assumptions in $M$ --- which other candidate literals are also forced.
This costs only a handful of unsatisfiable SAT calls, all of which
finish well inside the conflict cap.

When several MUSes of the same minimum size are available for the
same step, we then keep the one that justifies the most deductions.
This reduces the number of explanation steps without increasing the
size of any individual step's explanation, and is independent of the
choice of MUS strategy.

\paragraph{Aggressive conflict cap with auto-ramp.}

A small minority of SAT calls during MUS search hit pathological
constraint subsets that take seconds (occasionally minutes) to
resolve.  These tail calls dominated wall-clock time on harder
instances.  We address this by capping every SAT call at 100
conflicts and treating the limit hit as ``inconclusive'' rather than
as ``unsatisfiable'': inconclusive calls cause the surrounding
deletion-based search to keep the constraint conservatively, exactly
as if it were known to be required.

To recover correctness on instances where 100 conflicts is genuinely
too tight, we monitor the rolling fraction of inconclusive calls and
ramp the cap up by~$10\times$ whenever it exceeds a threshold.  Across
the 1\,960 \textsc{mus} and \textsc{core} runs in this benchmark,
ramping triggered on only five instance/method pairs, all going from
100 to 1\,000 and never further:

\begin{table}[h]
\centering
\small
\begin{tabular}{@{}lcc@{}}
\toprule
Strategy & Instances ramped & Total ramp events \\
\midrule
\textsc{mus}   & 2 & 2 \\
\textsc{core}  & 3 & 5 \\
\midrule
Total          & 5 / 980 & 7 \\
\bottomrule
\end{tabular}
\end{table}

\paragraph{Effect of tightening the cap.}

The previous default was a per-call conflict limit of 1\,000\,000 ---
effectively no limit, since SAT calls rarely approached it; the limit
existed only as a safety net against runaway searches.  Tightening it
to~100 reduced total wall-clock time across the suite to $0.70\times$
of the old runtime under \textsc{mus} and $0.90\times$ under
\textsc{core} (which already finished most calls in a handful of
conflicts).  The benefit concentrates on the families where
pathological tail calls were most common: Futoshiki~9$\times$9 hard
drops to $0.61\times$ of its previous runtime under \textsc{mus}, and
Tents~OA~10$\times$10 hard to $0.29\times$.

The cap is not free: discarding inconclusive calls means we
occasionally fail to remove a redundant constraint that a longer
search would have eliminated.  In aggregate this shows up as a small
shift in mean MUS size --- under the old default, mean MUS size
across all instances was 1.338 for \textsc{core} and 1.217 for
\textsc{mus}; under the new default, 1.333 and 1.211 respectively.
The change is within the noise expected from random tie-breaking
during MUS search, and the maximum-MUS-size column of
Table~\ref{tab:mus-comparison} is essentially unchanged.


% =============================================================
%  RESULTS DISCUSSION
% =============================================================

\subsection{Results}

Table~\ref{tab:mus-comparison} shows solve time and maximum MUS size
for each strategy.

\paragraph{Core is fast but imprecise.}

\textsc{core} is consistently the fastest strategy, often dramatically
so.  On hard instances where MUS search is expensive, the speed-up is
large: Futoshiki~9$\times$9 hard takes 9.52\,s with \textsc{mus} but
only 1.37\,s with \textsc{core} (6.9$\times$ faster), and
Sudoku~evil takes 5.99\,s vs.\ 1.15\,s (5.2$\times$).  The
only-assign variants show similar gaps: Futoshiki~OA 9$\times$9 hard
takes 4.00\,s with \textsc{mus} vs.\ 0.54\,s with \textsc{core}
(7.4$\times$), and Tents~OA 10$\times$10 hard takes 9.63\,s vs.\
2.18\,s (4.4$\times$).

The cost of this speed is explanation quality.  SAT solver cores are
not minimal and can be substantially larger than true MUSes.  On
Futoshiki~OA 9$\times$9 hard, \textsc{core} returns explanations with
a mean maximum size of~60.6 constraints, compared to 25.9 for
\textsc{mus} --- more than double.  Similar inflation appears across
all hard difficulty levels.  On easy instances where most MUSes are
already small (size~1--2), the difference is modest.

\paragraph{Tents is an outlier.}

Tents (without only-assign) shows no speed advantage for \textsc{core}
at all: 10$\times$10 instances take 25--26\,s regardless of strategy.
This suggests that for Tents, the bottleneck is not MUS search but
some other aspect of the solve loop (likely the large number of forced
deductions per step rather than the difficulty of justifying individual
deductions).

\paragraph{Summary.}

\textsc{core} offers a substantial speed-up (typically 2--7$\times$
on hard instances) at the cost of larger explanations
(1.4--3.4$\times$ more constraints).  For applications where
explanation size matters, \textsc{mus} remains the best strategy;
where speed matters and larger explanations are acceptable,
\textsc{core} is the clear winner.


% =============================================================
%  CASE STUDY: MIRACLE SUDOKU
% =============================================================

\subsection{Case study: the Miracle Sudoku}

The ``Miracle Sudoku'', popularised by Mitchell Lee on Cracking the
Cryptic, is a 9$\times$9 grid with only two given digits (a~1 in
row~5 column~3, a~2 in row~6 column~7) plus three extra restrictions
on top of standard Sudoku: orthogonally adjacent cells cannot contain
consecutive digits, no two cells a king's move apart share a digit,
and no two cells a knight's move apart share a digit.  Despite its
sparse clues, the puzzle has a unique solution.  It is notorious for
resisting every standard Sudoku technique --- the chains of
anti-knight and anti-king interactions are essential to any
explanation.

We ran the puzzle through Demystify with both strategies; the
contrast between them is sharper than on any benchmark instance.

\textsc{core} returns an explanation in \textbf{4.48\,s} of solve
time (10.9\,s wall-clock including Conjure compilation and parse
setup), built from 100 steps and 398 returned cores across 37\,033
SAT calls.  382 of the 398 cores (96\%) are unit cores, but the
remaining sixteen include four at sizes 15, 19, 20, and~38 ---
unminimised cores that bundle together far more constraints than are
strictly necessary.

\textsc{mus} takes \textbf{107\,s} of solve time (~25$\times$
longer), but produces a noticeably better explanation: 79 steps
(21~fewer than \textsc{core}) with 410 minimised MUSes whose
\emph{maximum} size is~4, and whose mean is 1.17.  Every individual
MUS in the explanation is therefore small enough to read at a glance,
where \textsc{core}'s size-38 cores would require careful unpacking.

The Miracle Sudoku is hard for two unrelated reasons.  Globally, only
two givens means narrow forced-move chains in the early steps and
hence many deductions before the explanation gathers momentum.
Locally, the cross-cutting anti-knight, anti-king and
non-consecutive constraints generate dense conflict graphs in which
SAT cores can be much larger than the underlying minimal explanation
--- exactly the regime where minimisation pays off.

Unlike every other puzzle in the corpus, the Miracle Sudoku is also
capable of producing genuinely hard \emph{individual} SAT calls.  In
an earlier run with the conflict cap disabled, 645 SAT calls each
took longer than two seconds, with the longest single call observed
at 42.4\,s and 87\,102 conflicts; 479 calls exceeded 10\,000
conflicts each.  Total solve time in that uncapped run was
4\,263\,s, against 107\,s under \textsc{mus} with the cap at~100 ---
a $40\times$ difference attributable almost entirely to the right
tail of the SAT-call distribution.  This is the case where the cap
earns its keep most clearly: the cap prevents the search from
chasing minimality through the worst chains, at the cost (on this
instance) of nothing visible in the resulting explanation.
